% Nouveau EGA V (cote25): archive cote25, site124 (PDF page125).
% English translation of the checked French annotated master.
% AI-translated by GPT-6 Astra (OpenAI, ChatGPT Pro mode).
\scanpage{25/124}
\mspage{27}

condition is expressed simply by requiring that for every $t\hand{\in} E(s)$,
the corresponding \unsure{section} u(s) of $\underline{L}_s$ on $X_s$ be \unsure{regular}, and that the
analogous condition hold after every extension of the base field.
One sees as \var{\over{\unsure{usual}}{\unsure{usual}}}{\unsure{usual}}{writing over}{ega25_25_124-v0001} that it suffices to test this condition on an
algebraic closure of k. To summarize:
\uline{Proposition} \var{\over{5}{4}}{4}{writing over}{ega25_25_124-v0002}.1. Le\hand{t} $g:X\hand{\to} S$ be a flat morphism loc of finite
presentation, $P=P(E)$ a projective bundle over X defined by $\underline{E}$ loc free of finite
type\var{\inserted{everywhere $\neq 0$, i.e. $P$ with nonempty fibers}}{ everywhere $\neq 0$, i.e. $P$ with nonempty fibers}{insertion}{ega25_25_124-v0004}, $\overset{\hand{\vee}}{P} = P(\overset{\hand{\vee}}{E})$ \hand{,} Then there is a bijection between
the set of linear systems of divisors on X/S parametrized by $\overset{\hand{\vee}}{P}$ ,
and the set of pairs(up to isomorphism) $(\underline{L},u)$, where $\underline{L}$ is an
invertible Module on X and u: $\underline{E}\hand{\to} g_{\hand{*}}(\underline{L})$ a homomorphism, such that
for every $s\hand{\in} S$ and every point t of $E(s)\otimes_{k(s)}k$ (for an arbitrary
\var{\margin{$\cdot$}}{\placedmargin{$\cdot$}}{marginal note placed}{ega25_25_124-v0005} extension\var{\inserted{$k$}}{ $k$}{insertion}{ega25_25_124-v0006} of k\unsure{\var{\over{(s)}{\illegible{1}}}{\illegible{1}}{writing over}{ega25_25_124-v0007}}, which one can suppose to be the algebraic closure of k(s)),
the corresponding section u(t) of $\underline{L}_{s\,k}$ on $X_{s\,k}$ is regular.

One will note that when the fibers of X are \var{\xout{($S_1$) and geom irred}}{}{typewriter strike-out}{ega25_25_124-v0008}
geometrically integral, this condition on u simply means that
for every $s\hand{\in} S$, $u(s): E(s)\hand{\rightarrow} H^{0}(X_s,\underline{L}_s)$ is injective, a fact which
should also be made explicit in \var{\over{5}{4}}{4}{writing over}{ega25_25_124-v0009}.1. ; \var{\deleted{\unsure{and|or}}}{}{deletion}{ega25_25_124-v0010} one should also recall
(for convenience of reference) the construction of the divisor D in
terms of $(\underline{L},u)$ as the divisor of the obvious section \var{\inserted{$\phi$}}{$\phi$}{insertion}{ega25_25_124-v0011} of $\underline{L}\otimes_{\underline{O}_S}\underline{O}_{\overset{\hand{\vee}}{P}}(1)$
defined by u.
\uline{Corollary} \var{\over{5}{4}}{4}{writing over}{ega25_25_124-v0012}.2. \var{\xout{Let L}}{}{typewriter strike-out}{ega25_25_124-v0013} Suppose $f: X\hand{\longrightarrow} S$ proper, flat, of finite
presentation and with geometrically integral fibers, \var{\xout{and}}{}{typewriter strike-out}{ega25_25_124-v0014} let $\underline{L}$ be an invertible
Module on X, and $P=P(\underline{E})$ a projective bundle over S as in 5.1.
\var{\xout{Finally let}}{}{typewriter strike-out}{ega25_25_124-v0015} There exists a Module Q of finite presentation on S and an
isomorphism of functors in the quasi-coherent $\underline{O}_S$-Module F :
\[ \mathrm{Hom}\,(Q,F) \hand{\xrightarrow{\ \sim\ }} g_{\hand{*}}(\underline{L}\otimes_{\underline{O}_S}F) \quad . \]
With this understood, the \var{\xout{families}}{}{typewriter strike-out}{ega25_25_124-v0016} linear systems of divisors on X parametrized
\tnote{25-124-1}{The first two doubts are in French ``section'' (the initial se) and ``régulière'' (initial ŕ/é). The overprint within ``habitude'' has before-candidates ah/ha and after-candidates ha/h; both are represented by doubtful ``usual'', with the overprint retained. The canceled doubtful \unsure{a} within ``Soit'' has no separate English letter; its handwritten final t is retained in English ``Let''. ``Loc'', ``geom'', and ``irred'' mean locally, geometrically, and irreducible. The initial u(s), the projective bundle said to be over X, and the later reference5.1 remain as in the source.}

\mspage{124}
